\chapter{Transactions and State Transitions}

\section{Purpose}

Chapter 1 gave us the state-machine model. Chapter 2 taught us to
identify assets and authority. Chapter 3 studies the object that usually
carries authority into the state machine:

\begin{quote}
A transaction is a signed, metered, ordered request to execute one or
more state transitions.
\end{quote}

That definition is intentionally more careful than ``a transaction sends
money.'' Some transactions transfer native assets. Some call contracts.
Some deploy code. Some carry multiple instructions. Some vote. Some
update validator sets. Some submit proofs. Some relay cross-chain
messages. Some fail but still consume fees. Some leave no application
state change but still affect liveness, ordering, fees, or user
expectations.

Transactions are where user intent, cryptographic authorization,
resource pricing, chain ordering, execution semantics, and application
logic meet. That makes them one of the richest security surfaces in
blockchain systems.

The central discipline of this chapter is:

\begin{lstlisting}
Do not ask only what a transaction says.
Ask what state it can change, under whose authority, at what cost,
in which order, with what failure behavior, and with which observable result.
\end{lstlisting}

\section{Learning Objectives}

By the end of this chapter, you should be able to:

\begin{itemize}
\tightlist
\item
  Describe transaction anatomy across UTXO, account, instruction-based,
  object-based, and Cosmos-style systems.
\item
  Explain how signatures, nonces, sequence numbers, recent blockhashes,
  object versions, and chain IDs limit replay.
\item
  Explain fees, gas, compute budgets, fee markets, and resource
  exhaustion as security mechanisms, not just UX concerns.
\item
  Interpret transaction results, receipts, logs, events, return data,
  status codes, and reverts without confusing them for application
  truth.
\item
  Distinguish syntactic validity, mempool acceptance, block inclusion,
  execution success, finality, and protocol safety.
\item
  Identify valid transactions that can still produce unsafe states.
\end{itemize}

\section{The Core Model}

A transaction has three layers:

{\def\LTcaptype{none} % do not increment counter
\begin{longtable}[]{@{}
  >{\raggedright\arraybackslash}p{(\linewidth - 4\tabcolsep) * \real{0.3333}}
  >{\raggedright\arraybackslash}p{(\linewidth - 4\tabcolsep) * \real{0.3333}}
  >{\raggedright\arraybackslash}p{(\linewidth - 4\tabcolsep) * \real{0.3333}}@{}}
\toprule\noalign{}
\begin{minipage}[b]{\linewidth}\raggedright
Layer
\end{minipage} & \begin{minipage}[b]{\linewidth}\raggedright
Review question
\end{minipage} & \begin{minipage}[b]{\linewidth}\raggedright
Typical fields or evidence
\end{minipage} \\
\midrule\noalign{}
\endhead
\bottomrule\noalign{}
\endlastfoot
Authorization & Who is asking for this transition, and how is authority
proven? & Signatures, signers, account keys, witnesses, approvals,
roles \\
Execution request & What operation, message, instruction, script, or
call should run? & Calldata, script, instruction data, message type,
target program/contract \\
Resource and ordering envelope & What determines whether, when, and how
the request can be included? & Fees, gas, compute budget, nonce,
sequence, blockhash, expiration, chain ID \\
\end{longtable}
}

Different chains package these layers differently.

Ethereum documentation describes transactions as cryptographically
signed instructions from accounts that update network state.\footnote{ethereum.org,
  ``Transactions'', https://ethereum.org/developers/docs/transactions/.}
Bitcoin's developer guide describes transactions as composed of inputs
and outputs, where each input spends a previous output and each output
waits as a UTXO until later spent.\footnote{Bitcoin Developer Guide,
  ``Transactions'',
  https://developer.bitcoin.org/devguide/transactions.html.} Solana
documentation describes transactions as containing one or more
instructions, account signatures, and a recent blockhash; all
instructions in a transaction are processed together and failure of one
instruction fails the entire transaction.\footnote{Solana Foundation,
  ``Transactions'', https://solana.com/docs/core/transactions.} Cosmos
SDK documentation describes a lifecycle that moves from creation and
signing, to mempool checks, to block inclusion, to deterministic
execution and committed state changes.\footnote{Cosmos SDK v0.47,
  ``Transaction Lifecycle'',
  https://docs.cosmos.network/sdk/v0.47/learn/beginner/tx-lifecycle.}

Those systems look different, but they all answer the same questions:

\begin{lstlisting}
Who authorized the input?
What state does it read?
What state does it write?
What resources does it consume?
How is replay prevented?
What happens if it fails?
When is the result accepted?
\end{lstlisting}

This chapter is about those questions.

The lifecycle also matters. A transaction can be safe at one point and
unsafe at another:

\begin{lstlisting}
flowchart LR
    A[Constructed] --> B[Signed]
    B --> C[Submitted]
    C --> D[Propagated]
    D --> E[Pending]
    E --> F[Included]
    F --> G[Executed]
    G --> H[Confirmed]
    H --> I[Finalized / settled]
\end{lstlisting}

\section{Transaction Anatomy Across Chain Models}

\subsection{Transaction Anatomy Is a Security Interface}

A transaction is not only a data structure. It is the boundary between
intent and state transition.

The fields inside a transaction determine:

\begin{itemize}
\tightlist
\item
  Who authorizes it.
\item
  Which state objects or accounts it can touch.
\item
  Which code path runs.
\item
  How much value moves.
\item
  How much resource the sender is willing to pay for.
\item
  Whether the transaction can be replayed.
\item
  Whether it expires.
\item
  Whether it can be replaced.
\item
  Whether it is valid only on a particular chain or domain.
\item
  What data validators, wallets, relayers, builders, and applications
  can inspect.
\end{itemize}

This means transaction anatomy is security-critical. A missing field,
ambiguous field, incorrectly signed field, or misunderstood field can
become an exploit primitive.

\subsection{UTXO Transactions: Consume Outputs, Create Outputs}

Bitcoin-style transactions are built from inputs and outputs.

Each input references a previous output by transaction ID and output
index. Each output assigns value to a spending condition. A wallet
balance is a derived view over spendable unspent outputs, not a single
mutable account balance.\footnote{Bitcoin Developer Guide,
  ``Transactions'',
  https://developer.bitcoin.org/devguide/transactions.html.}

Abstractly:

\begin{lstlisting}
Transaction:
  version
  inputs:
    previous_txid
    previous_output_index
    unlocking_data / witness
    sequence
  outputs:
    amount
    locking_script
  locktime
\end{lstlisting}

The state transition is:

\begin{lstlisting}
remove consumed UTXOs
add newly created UTXOs
\end{lstlisting}

Validity depends on conditions such as:

\begin{itemize}
\tightlist
\item
  Referenced outputs exist and are unspent.
\item
  Unlocking data satisfies each referenced output's spending condition.
\item
  Input value is sufficient to cover output value.
\item
  No output is spent twice.
\item
  Locktime and sequence constraints are satisfied.
\item
  Consensus and standardness rules are met.
\end{itemize}

Security consequences:

\begin{itemize}
\tightlist
\item
  There is no account nonce in the Ethereum sense.
\item
  Replay is constrained because spending consumes specific outputs.
\item
  Change outputs are normal outputs and can leak privacy or be
  mishandled.
\item
  Fee is typically implicit: input value minus output value.
\item
  Transaction malleability historically mattered because downstream
  transactions may reference transaction IDs.
\item
  Script conditions define spending authority.
\end{itemize}

UTXO analysis starts with:

\begin{lstlisting}
Which outputs are consumed?
Which new outputs are created?
Which spending conditions were satisfied?
Which value disappeared as fee?
Which future spending paths now exist?
\end{lstlisting}

\subsection{Account Transactions: Mutate Account and Contract State}

Account-based systems represent state differently. Ethereum transactions
are initiated by externally owned accounts and can transfer ETH, deploy
contracts, or call contract code.\footnote{ethereum.org,
  ``Transactions'', https://ethereum.org/developers/docs/transactions/.}

A modern Ethereum transaction may include:

\begin{lstlisting}
from
to
nonce
value
input data / calldata
gas limit
max fee per gas
max priority fee per gas
chain-specific signature data
access list or transaction type fields
signature
\end{lstlisting}

Ethereum.org lists these core fields, including sender, recipient,
signature, nonce, value, input data, gas limit, max priority fee, and
max fee.\footnote{ethereum.org, ``Transactions'',
  https://ethereum.org/developers/docs/transactions/.}

The state transition may update:

\begin{itemize}
\tightlist
\item
  Sender nonce.
\item
  Native ETH balances.
\item
  Contract storage.
\item
  Account code if a contract is created.
\item
  Logs emitted during execution.
\item
  Gas accounting and fee distribution.
\end{itemize}

In the simplest ETH transfer:

\begin{lstlisting}
balances[sender] -= value + fee
balances[recipient] += value
nonce[sender] += 1
base_fee portion is burned
priority fee goes to block proposer
\end{lstlisting}

In a contract call:

\begin{lstlisting}
nonce[sender] += 1
call target with calldata and value
execute EVM code
update storage, balances, logs, and gas accounting according to execution
\end{lstlisting}

Security consequences:

\begin{itemize}
\tightlist
\item
  A single \passthrough{\lstinline!to!} address can hide complex logic
  if it is a contract.
\item
  Calldata may be opaque to users and wallets.
\item
  A transaction that sends zero native value may still move tokens
  through contract logic.
\item
  Nonces serialize transactions from an account.
\item
  Failed execution can still consume fees.
\item
  Logs may tell applications what happened, but logs are not the same as
  state.
\end{itemize}

Account-model analysis starts with:

\begin{lstlisting}
Who signed?
Which nonce is consumed?
Which target is called?
What calldata is supplied?
What value is attached?
Which storage, balances, logs, and subcalls can change?
What happens on revert?
\end{lstlisting}

\subsection{Instruction-Based Transactions: Multiple Operations, Shared
Atomicity}

Solana uses transactions containing one or more instructions. The
transaction includes signatures from accounts that authorize changes and
a recent blockhash. Solana's docs emphasize that instructions in a
transaction are processed together: if any instruction fails, the entire
transaction fails and state changes are reverted; fees are still charged
on failure.\footnote{Solana Foundation, ``Transactions'',
  https://solana.com/docs/core/transactions.}

A Solana transaction contains, at a high level:

\begin{lstlisting}
signatures
message:
  header
  account addresses
  recent blockhash
  compiled instructions:
    program id
    account indexes
    instruction data
\end{lstlisting}

Security consequences:

\begin{itemize}
\tightlist
\item
  Atomicity spans multiple instructions.
\item
  The caller supplies the account list.
\item
  Programs must verify that supplied accounts are the expected accounts.
\item
  Signer and writable flags are part of the authority boundary.
\item
  A recent blockhash limits transaction lifetime.
\item
  Transaction size and compute limits are first-class constraints.
\item
  Cross-program invocations can propagate privileges in ways reviewers
  must understand.
\end{itemize}

Solana analysis starts with:

\begin{lstlisting}
Which instructions execute?
Which accounts are passed to each instruction?
Which accounts signed?
Which accounts are writable?
Which program owns each account?
Which account constraints are enforced by the program?
Which instruction fails if assumptions are false?
\end{lstlisting}

\subsection{Object-Based Transactions: Inputs, Versions, and Effects}

Object-centric systems such as Sui represent many assets and state
components as objects. Sui documentation describes every update on Sui
as happening through a transaction and distinguishes user-submitted
programmable transaction blocks from system transactions submitted by
validators.\footnote{Sui Foundation, ``Transaction Overview'',
  https://docs.sui.io/develop/transactions/txn-overview.}

Sui transaction metadata includes fields such as sender address, gas
input object, gas price, gas budget, epoch, transaction type,
authenticator, and optional expiration.\footnote{Sui Foundation,
  ``Transaction Overview'',
  https://docs.sui.io/develop/transactions/txn-overview.}

The object model changes security analysis:

\begin{itemize}
\tightlist
\item
  Objects are explicit inputs and outputs.
\item
  Ownership and object versions matter.
\item
  Parallel execution can be possible when transactions touch disjoint
  objects.
\item
  Shared objects introduce contention and ordering concerns.
\item
  Gas is paid through a gas object.
\item
  Expiration can bound when a transaction may execute.
\end{itemize}

Sui-style analysis starts with:

\begin{lstlisting}
Which objects are inputs?
Which objects are mutated, created, transferred, wrapped, or destroyed?
Which object versions are assumed?
Which objects are owned, shared, or immutable?
What is the gas object?
What are the effects?
\end{lstlisting}

\subsection{Cosmos-Style Transactions: Messages Through an Application
Pipeline}

Cosmos SDK transactions usually contain one or more messages. The
transaction passes through creation, signing, mempool checks, block
proposal, execution, and commit.

The Cosmos SDK transaction lifecycle distinguishes:

\begin{itemize}
\tightlist
\item
  Transaction creation and signing.
\item
  Gas and fee settings.
\item
  \passthrough{\lstinline!CheckTx!} before mempool acceptance.
\item
  Stateless and stateful checks.
\item
  An ante handler for signatures, sequence numbers, fee deduction, and
  related operations.
\item
  \passthrough{\lstinline!DeliverTx!} during block execution.
\item
  Commit of state changes after consensus.\footnote{Cosmos SDK v0.47,
    ``Transaction Lifecycle'',
    https://docs.cosmos.network/sdk/v0.47/learn/beginner/tx-lifecycle.}
\end{itemize}

Security consequences:

\begin{itemize}
\tightlist
\item
  A transaction can pass initial checks and still fail later.
\item
  Local mempool fee policies can differ from deterministic block
  execution rules.
\item
  Sequence numbers and fees help prevent replay and spam.
\item
  Multiple messages may execute under one transaction.
\item
  Application modules define message handlers that cause state
  transitions.
\end{itemize}

Cosmos-style analysis starts with:

\begin{lstlisting}
Which messages are inside the transaction?
Which module handles each message?
Which ante-handler checks occur?
Which account sequence is consumed?
Which fees are deducted?
Which checks happen in mempool only, and which happen during block execution?
\end{lstlisting}

\subsection{Cross-Model Comparison}

{\def\LTcaptype{none} % do not increment counter
\begin{longtable}[]{@{}
  >{\raggedright\arraybackslash}p{(\linewidth - 8\tabcolsep) * \real{0.2000}}
  >{\raggedright\arraybackslash}p{(\linewidth - 8\tabcolsep) * \real{0.2000}}
  >{\raggedright\arraybackslash}p{(\linewidth - 8\tabcolsep) * \real{0.2000}}
  >{\raggedright\arraybackslash}p{(\linewidth - 8\tabcolsep) * \real{0.2000}}
  >{\raggedright\arraybackslash}p{(\linewidth - 8\tabcolsep) * \real{0.2000}}@{}}
\toprule\noalign{}
\begin{minipage}[b]{\linewidth}\raggedright
Model
\end{minipage} & \begin{minipage}[b]{\linewidth}\raggedright
Primary input shape
\end{minipage} & \begin{minipage}[b]{\linewidth}\raggedright
Replay boundary
\end{minipage} & \begin{minipage}[b]{\linewidth}\raggedright
Resource boundary
\end{minipage} & \begin{minipage}[b]{\linewidth}\raggedright
Common security focus
\end{minipage} \\
\midrule\noalign{}
\endhead
\bottomrule\noalign{}
\endlastfoot
UTXO & Inputs consume outputs; outputs create new claims & Spent outputs
& Fee from input-output difference, size/weight & Double spend, script
conditions, change, pinning \\
Ethereum account & Signed account transaction with nonce, gas, calldata
& Account nonce and chain/domain separation & Gas limit and fee fields &
Calldata intent, nonce ordering, reverts, gas griefing \\
Solana & Signed transaction with account list, instructions, recent
blockhash & Recent blockhash, durable nonce, account/signature checks &
Compute units, size limits, fees & Account substitution, signer/writable
flags, CPI \\
Sui/object & Transaction block over objects and gas object & Object
versions, sender/authenticator, epoch/expiration & Gas budget, object
limits & Object ownership, shared objects, effects \\
Cosmos SDK & Signed transaction with messages, fees, sequence & Account
sequence, mempool cache, timeout height & Gas wanted, fees, min gas
price & Ante handler, module message routing, deterministic deliver \\
\end{longtable}
}

No single anatomy is ``the blockchain transaction model.'' The right
security model comes from the state representation and execution rules
of the chain being reviewed.

For chapter-level review, keep the anatomy table compact:

{\def\LTcaptype{none} % do not increment counter
\begin{longtable}[]{@{}
  >{\raggedright\arraybackslash}p{(\linewidth - 6\tabcolsep) * \real{0.2500}}
  >{\raggedright\arraybackslash}p{(\linewidth - 6\tabcolsep) * \real{0.2500}}
  >{\raggedright\arraybackslash}p{(\linewidth - 6\tabcolsep) * \real{0.2500}}
  >{\raggedright\arraybackslash}p{(\linewidth - 6\tabcolsep) * \real{0.2500}}@{}}
\toprule\noalign{}
\begin{minipage}[b]{\linewidth}\raggedright
Model
\end{minipage} & \begin{minipage}[b]{\linewidth}\raggedright
Fields the reviewer must locate
\end{minipage} & \begin{minipage}[b]{\linewidth}\raggedright
What proves authority?
\end{minipage} & \begin{minipage}[b]{\linewidth}\raggedright
What shows success?
\end{minipage} \\
\midrule\noalign{}
\endhead
\bottomrule\noalign{}
\endlastfoot
Bitcoin/UTXO & Inputs, previous outputs, witnesses/scripts, output
amounts, locktime, sequence & Satisfaction of each consumed output's
spending condition & Transaction remains in canonical history; intended
outputs exist and inputs are spent \\
Ethereum/account & \passthrough{\lstinline!from!},
\passthrough{\lstinline!to!}, nonce, value, calldata, gas/fee fields,
chain ID, signature & Sender signature, contract authorization checks,
approvals or permits & Receipt status, state changes, logs as evidence,
sufficient settlement \\
Solana/instruction & Signatures, account list, writable/signer flags,
recent blockhash, instructions & Required signatures and
program-enforced account constraints & Transaction confirmation, account
state changes, emitted logs if relevant \\
Object/resource-based & Sender/authenticator, input objects/resources,
versions, gas object, effects, expiration & Ownership, capabilities,
authenticators, module rules & Object effects, version changes, event
evidence, settlement status \\
\end{longtable}
}

This table is deliberately shallow. Part III will handle each execution
model in depth. Here the purpose is to learn what must be inspected
before trusting a transaction outcome.

\section{Signatures, Nonces, Sequence Numbers, and Replay Prevention}

\subsection{Signatures Authorize a Specific Message, Not a Vague
Intention}

A signature does not prove that a user understood what they were doing.
It proves that some signing authority produced a valid cryptographic
signature over a particular message under a particular scheme.

This distinction matters.

A user may think they are signing:

\begin{lstlisting}
Swap 100 USDC for at least 0.04 ETH.
\end{lstlisting}

The wallet may actually present or submit:

\begin{lstlisting}
approve(spender = attacker, amount = unlimited)
\end{lstlisting}

The cryptographic signature can be valid while the human intention was
violated. Blockchain security must therefore separate:

\begin{lstlisting}
cryptographic authorization
from
informed consent
from
application safety
\end{lstlisting}

Chapter 4 will go deeper into keys and signatures. For this chapter,
remember:

\begin{lstlisting}
A transaction signature authorizes exactly the bytes and domain that the signature scheme verifies.
Everything not bound into the signed message is a potential ambiguity.
\end{lstlisting}

\subsection{Replay Is Reuse of Authority in the Wrong Context}

A replay attack occurs when a valid authorization is reused where it
should not be valid.

Replay can happen across:

\begin{itemize}
\tightlist
\item
  Time: the same transaction or message is accepted twice.
\item
  Chain: a transaction valid on one chain is valid on another.
\item
  Contract: a signature intended for one contract is accepted by
  another.
\item
  Function: a signature intended for one operation authorizes another.
\item
  User/account: a message is valid for an unintended signer or account.
\item
  Bridge domain: a cross-chain message is executed on the wrong
  destination or more than once.
\item
  Fork or clone: old signatures work after migration or redeployment.
\end{itemize}

Replay prevention is not one mechanism. It is a family of mechanisms
that bind authority to a specific context and consume it when used.

\subsection{Ethereum Nonces Serialize Account Transactions}

Ethereum transactions include a nonce: a sequential counter for
transactions from the sender account.\footnote{ethereum.org,
  ``Transactions'', https://ethereum.org/developers/docs/transactions/.}

The nonce serves several purposes:

\begin{itemize}
\tightlist
\item
  It makes two otherwise identical transactions distinct if their nonces
  differ.
\item
  It prevents the same signed transaction from being included twice on
  the same chain.
\item
  It orders transactions from one account.
\item
  It allows replacement behavior in pending transaction pools.
\end{itemize}

But nonces also introduce operational risks:

\begin{itemize}
\tightlist
\item
  A low-fee transaction can block later higher-nonce transactions from
  the same account.
\item
  A replacement transaction must satisfy node replacement rules.
\item
  A signing system that mismanages nonces can create stuck or
  conflicting transactions.
\item
  Parallel signing from the same account can produce nonce races.
\item
  Transaction cancellation is usually replacement with a same-nonce
  transaction, not erasure of history.
\end{itemize}

Security questions:

\begin{lstlisting}
Who assigns nonces?
Can two systems sign with the same nonce?
What happens if a transaction is stuck?
Can a malicious or compromised signer block future transactions?
Can an attacker force a protocol bot into nonce contention?
\end{lstlisting}

\subsection{Chain IDs and Domain Separation Prevent Cross-Chain Replay}

Ethereum's EIP-155 introduced simple replay protection by including the
chain ID in the transaction signing scheme after the Spurious Dragon
fork.\footnote{Vitalik Buterin, ``EIP-155: Simple Replay Attack
  Protection'', https://eips.ethereum.org/EIPS/eip-155.} This prevents a
transaction signed for one EIP-155-protected chain from being valid on
another chain with a different chain ID.

The security principle is broader than Ethereum:

\begin{lstlisting}
Every authorization should be bound to the domain where it is intended to be valid.
\end{lstlisting}

Domains may include:

\begin{itemize}
\tightlist
\item
  Chain ID.
\item
  Contract address.
\item
  Program ID.
\item
  Verifying contract.
\item
  Application name.
\item
  Version.
\item
  Source and destination chain.
\item
  Message nonce.
\item
  User account.
\item
  Expiration.
\end{itemize}

If a signature says only ``approve 100 tokens'' but not which chain,
which contract, which token, which spender, which deadline, and which
nonce, it may authorize more than intended.

Domain separation turns vague authority into scoped authority.

\subsection{UTXO Replay Prevention Comes From Consumption}

In a UTXO system, replay prevention is tied to spending specific
outputs.

Once an output is spent in the canonical chain, another transaction
trying to spend the same output conflicts. A transaction cannot be
accepted twice because its inputs would no longer be unspent.

But UTXO systems still have replay concerns:

\begin{itemize}
\tightlist
\item
  Transactions can be replayed across chains after forks if rules and
  UTXOs remain compatible.
\item
  Signatures may authorize different transaction variants depending on
  sighash modes.
\item
  Unconfirmed transaction replacement and pinning can affect propagation
  and settlement.
\item
  Child transactions may depend on parent transaction IDs or witness
  behavior.
\item
  Time locks and sequence values affect when spending is valid.
\end{itemize}

The main lesson is:

\begin{lstlisting}
UTXO replay prevention is based on consuming unique state, not incrementing an account counter.
\end{lstlisting}

\subsection{Solana Recent Blockhashes and Durable Nonces}

Solana transactions include a recent blockhash. Solana documentation
states that a transaction's recent blockhash is valid for 150
slots.\footnote{Solana Foundation, ``Transactions'',
  https://solana.com/docs/core/transactions.} This creates a
time-bounded replay window.

Security consequences:

\begin{itemize}
\tightlist
\item
  Old transactions expire and cannot be included forever.
\item
  Offline signing needs special handling.
\item
  Users and bots must manage expiration and resubmission.
\item
  A transaction can fail because it is too old, even if signatures and
  instruction data are otherwise correct.
\end{itemize}

Solana also supports durable nonces for offline signing workflows.
Durable nonces solve a real operational problem but introduce their own
lifecycle: nonce account creation, nonce advancement, authority, and
failure behavior.

Ask:

\begin{lstlisting}
Is replay bounded by blockhash expiry, durable nonce state, or both?
Who controls the nonce authority?
Can old signed transactions remain valid longer than expected?
Can a nonce account be advanced or consumed by an unexpected party?
\end{lstlisting}

\subsection{Cosmos Sequence Numbers and Timeout Heights}

Cosmos SDK documentation describes account sequence numbers as part of
replay prevention and notes that fees and sequence counters help
distinguish replayed transactions from identical but valid
transactions.\footnote{Cosmos SDK v0.47, ``Transaction Lifecycle'',
  https://docs.cosmos.network/sdk/v0.47/learn/beginner/tx-lifecycle.}

Cosmos transactions can also use timeout heights, limiting commitment
after a certain block height.\footnote{Cosmos SDK v0.47, ``Transaction
  Lifecycle'',
  https://docs.cosmos.network/sdk/v0.47/learn/beginner/tx-lifecycle.}

Security consequences:

\begin{itemize}
\tightlist
\item
  Sequence numbers prevent repeated acceptance of the same account
  authorization.
\item
  Timeout height limits stale transaction execution.
\item
  Fee checks affect spam resistance and mempool admission.
\item
  A transaction may pass local checks but later fail during delivery.
\end{itemize}

Sequence numbers are not merely bookkeeping. They are authority
consumption.

\subsection{Object Versions and Expiration}

Object-based systems can use object versions, owned object references,
epochs, and expiration rules to prevent stale or conflicting
transitions.

Sui documentation describes a transaction flow where objects are inputs
and outputs, and transaction metadata includes epoch and optional
expiration.\footnote{Sui Foundation, ``Transaction Overview'',
  https://docs.sui.io/develop/transactions/txn-overview.}

Security consequences:

\begin{itemize}
\tightlist
\item
  A transaction that uses an old object version may no longer be valid.
\item
  Shared objects require consensus and ordering.
\item
  Owned objects may allow parallel execution when independent.
\item
  Expiration prevents delayed execution after context changes.
\end{itemize}

Object models make some conflicts explicit:

\begin{lstlisting}
This transaction is not just from Alice.
It consumes or mutates these specific object versions.
\end{lstlisting}

\subsection{Replay Protection Must Match the Asset}

Replay protection is only correct if it protects the asset or authority
being consumed.

Examples:

\begin{itemize}
\tightlist
\item
  A transaction nonce protects Ethereum account-level transaction
  replay.
\item
  A permit nonce protects token approval signature replay.
\item
  A bridge message nonce protects cross-chain message replay.
\item
  A withdrawal nullifier protects double-claiming.
\item
  A UTXO input protects double-spending a specific output.
\item
  A governance proposal ID prevents accidental duplicate execution only
  if execution status is tracked.
\item
  A rollup output root claim needs finality, proof, and claim-status
  tracking.
\end{itemize}

Do not say ``we have a nonce'' as if that settles the matter.

Ask:

\begin{lstlisting}
What authority is the nonce consuming?
Is the nonce scoped to the correct signer, chain, contract, and operation?
Can the nonce be skipped, reused, reset, or shared across domains?
Is there an expiration?
Can the same effect occur through another path?
\end{lstlisting}

\section{Fees, Gas, Compute Budgets, and Resource Limits}

\subsection{Resource Pricing Is Part of Consensus Security}

Blockchains cannot allow unbounded computation, storage growth,
signature verification, or network bandwidth consumption. Public systems
need resource limits because any participant can submit inputs.

Fees and metering serve several purposes:

\begin{itemize}
\tightlist
\item
  Prevent unlimited computation.
\item
  Allocate scarce block space.
\item
  Compensate or reward block producers.
\item
  Discourage spam.
\item
  Price storage, execution, and signature verification.
\item
  Create ordering incentives.
\item
  Bound denial-of-service impact.
\end{itemize}

Fees are not only economics. They are a security control.

If execution were free and unbounded, attackers could force every node
to do arbitrary work. A blockchain cannot rely on ``users will not
submit expensive nonsense.'' Permissionless execution means attackers
will submit whatever the rules allow if it benefits them or harms
others.

\subsection{Ethereum Gas}

Ethereum gas measures the computational and storage work required to
execute transactions. Ethereum.org describes gas as the fee required to
conduct transactions or execute contracts, with each transaction
consuming gas units and fees depending on gas used and price per
gas.\footnote{ethereum.org, ``Gas and Fees'',
  https://ethereum.org/developers/docs/gas/.}

The EIP-1559 transaction fee market introduced a base fee that is burned
and a priority fee paid to validators, with transactions specifying max
fee and max priority fee.\footnote{Vitalik Buterin et al., ``EIP-1559:
  Fee Market Change for ETH 1.0 Chain'',
  https://eips.ethereum.org/EIPS/eip-1559.} Ethereum.org's transaction
page describes the \passthrough{\lstinline!gasLimit!},
\passthrough{\lstinline!maxPriorityFeePerGas!}, and
\passthrough{\lstinline!maxFeePerGas!} fields.\footnote{ethereum.org,
  ``Transactions'', https://ethereum.org/developers/docs/transactions/.}

Security consequences:

\begin{itemize}
\tightlist
\item
  If the gas limit is too low, execution can fail.
\item
  If execution runs out of gas, state changes revert but gas is consumed
  according to protocol rules.
\item
  If fees are too low, the transaction may remain pending or be ignored.
\item
  Complex loops over user-controlled state can become denial-of-service
  vulnerabilities.
\item
  Attackers can grief systems by forcing expensive paths, even when they
  cannot steal funds.
\item
  Protocols that assume an operation will always fit within block gas
  limits can fail as state grows.
\item
  Gas costs can change across forks or differ across chains, affecting
  brittle assumptions.
\end{itemize}

Gas analysis asks:

\begin{lstlisting}
Can this transition become too expensive to execute?
Can an attacker increase the cost for other users?
Can a required recovery or withdrawal path run out of gas?
Does failure consume fees without producing the intended state?
Can a low-level call fail because of gas forwarding assumptions?
\end{lstlisting}

\subsection{Solana Compute Budgets and Fees}

Solana charges a base fee per signature and optional prioritization
fees. Solana's fee documentation describes a base fee and an optional
prioritization fee, with the prioritization fee increasing the
likelihood that the current leader schedules the transaction ahead of
competitors.\footnote{Solana Foundation, ``Fees'',
  https://solana.com/docs/core/fees.}

Solana also has compute unit limits. Its documentation lists default and
maximum compute-unit limits and transaction size limits.\footnote{Solana
  Foundation, ``Fees'', https://solana.com/docs/core/fees.} The
transaction page lists a 1,232-byte size limit, 64-byte Ed25519
signatures, and blockhash expiry.\footnote{Solana Foundation,
  ``Transactions'', https://solana.com/docs/core/transactions.}

Security consequences:

\begin{itemize}
\tightlist
\item
  Transaction size limits constrain how many accounts and instructions
  can be included.
\item
  Compute limits constrain program complexity.
\item
  Account locking and writable accounts affect parallelism and
  contention.
\item
  Prioritization fees affect scheduling.
\item
  Failed transactions can still pay fees.
\item
  Bots can create congestion through failed or competing transactions.
\end{itemize}

Solana resource analysis asks:

\begin{lstlisting}
Can the transaction fit within size and account limits?
Can the program fit within compute limits under worst-case state?
Can attackers lock accounts or create contention?
Can they force expensive CPIs?
Can a required transaction fail because too many accounts are needed?
\end{lstlisting}

\subsection{Cosmos Gas and Local Mempool Fee Policy}

Cosmos SDK documentation distinguishes user-specified gas, gas prices,
fees, local minimum gas prices, mempool checks, and deterministic
delivery.\footnote{Cosmos SDK v0.47, ``Transaction Lifecycle'',
  https://docs.cosmos.network/sdk/v0.47/learn/beginner/tx-lifecycle.}

Important nuance: local mempool admission can depend on each node's
local \passthrough{\lstinline!min-gas-prices!}, but deterministic block
execution cannot depend on local settings. Cosmos documentation notes
that \passthrough{\lstinline!DeliverTx!} does not compare gas prices to
local \passthrough{\lstinline!min-gas-prices!} because differing local
values would yield nondeterministic results.\footnote{Cosmos SDK v0.47,
  ``Transaction Lifecycle'',
  https://docs.cosmos.network/sdk/v0.47/learn/beginner/tx-lifecycle.}

Security consequences:

\begin{itemize}
\tightlist
\item
  A transaction may be rejected by one node's mempool and accepted by
  another.
\item
  Fee policy affects propagation and inclusion.
\item
  Gas estimation can undercount worst-case execution.
\item
  Ante-handler logic affects signatures, sequences, and fee deduction.
\item
  Block gas limits can prevent large batches or emergency operations.
\end{itemize}

The security lesson is general:

\begin{lstlisting}
Mempool policy may be local.
Consensus execution must be deterministic.
Do not confuse the two.
\end{lstlisting}

\subsection{UTXO Fees and Transaction Pinning}

In Bitcoin-style systems, fees are commonly implicit: the difference
between total input value and total output value.

\begin{lstlisting}
fee = sum(inputs) - sum(outputs)
\end{lstlisting}

Fee policy affects:

\begin{itemize}
\tightlist
\item
  Whether a transaction propagates.
\item
  Whether miners include it.
\item
  Whether it can be replaced.
\item
  Whether child-pays-for-parent can rescue it.
\item
  Whether an attacker can pin a transaction and prevent timely
  confirmation.
\end{itemize}

UTXO fee analysis asks:

\begin{lstlisting}
What fee does this transaction pay?
Can it be replaced?
Can a child transaction increase its effective fee?
Can an attacker create a low-fee dependency?
Can time-sensitive protocols fail if confirmation is delayed?
\end{lstlisting}

This will matter later for payment channels, bridges, vault exits, and
liquidation systems that depend on timely confirmation.

\subsection{Resource Limits Create Liveness Risks}

Resource limits do not only stop attackers. They can stop honest users.

Examples:

\begin{itemize}
\tightlist
\item
  A withdrawal loop over all users becomes impossible after the user set
  grows.
\item
  A liquidation path needs too much gas during market stress.
\item
  A bridge proof exceeds transaction size limits.
\item
  A governance proposal includes too many actions to execute atomically.
\item
  A Solana instruction needs more accounts than the transaction can
  carry.
\item
  A Cosmos module message consumes more gas than expected under extreme
  state.
\item
  A rollup escape hatch is too expensive during L1 congestion.
\end{itemize}

The invariant may still be true:

\begin{lstlisting}
Every user has a claim.
\end{lstlisting}

But liveness may fail:

\begin{lstlisting}
No user can practically execute the transition that realizes the claim.
\end{lstlisting}

Security reviews must test not only ``does the code work on an example''
but:

\begin{lstlisting}
Does the required transition remain executable under worst-case state,
market stress, fee spikes, and adversarial input?
\end{lstlisting}

\subsection{Fees Shape Adversarial Strategy}

Attackers are cost-sensitive. Fees determine which attacks are
profitable or feasible.

Examples:

\begin{itemize}
\tightlist
\item
  A spam attack may be uneconomic if fees are high enough.
\item
  A sandwich attack may be profitable only if expected extraction
  exceeds gas and priority fees.
\item
  A governance griefing attack may depend on cheap proposal creation.
\item
  A liquidation race may be won by paying higher priority fees.
\item
  A bridge relayer may stop relaying if fees exceed compensation.
\item
  A keeper network may fail if executing maintenance tasks is
  unprofitable.
\end{itemize}

Security analysis should include:

\begin{lstlisting}
What does the attacker pay?
What can the attacker gain?
Who bears failed transaction costs?
Who is reimbursed?
Who has incentives to include, relay, execute, or ignore the transaction?
\end{lstlisting}

Fees are part of the threat model.

\section{Execution Results: Receipts, Logs, Events, and Reverts}

\subsection{Inclusion Is Not Execution Success}

A transaction can exist in several states:

\begin{itemize}
\tightlist
\item
  Constructed but unsigned.
\item
  Signed but not broadcast.
\item
  Broadcast but not seen by most nodes.
\item
  Accepted into some mempools.
\item
  Included in a block.
\item
  Executed successfully.
\item
  Executed and reverted or failed.
\item
  Reorged out.
\item
  Finalized or settled.
\item
  Reflected correctly by an RPC, indexer, or UI.
\end{itemize}

These are different facts.

Do not say ``the transaction went through'' unless you specify which
fact you mean.

{\def\LTcaptype{none} % do not increment counter
\begin{longtable}[]{@{}
  >{\raggedright\arraybackslash}p{(\linewidth - 4\tabcolsep) * \real{0.3333}}
  >{\raggedright\arraybackslash}p{(\linewidth - 4\tabcolsep) * \real{0.3333}}
  >{\raggedright\arraybackslash}p{(\linewidth - 4\tabcolsep) * \real{0.3333}}@{}}
\toprule\noalign{}
\begin{minipage}[b]{\linewidth}\raggedright
Claim
\end{minipage} & \begin{minipage}[b]{\linewidth}\raggedright
What it means
\end{minipage} & \begin{minipage}[b]{\linewidth}\raggedright
What it does not prove
\end{minipage} \\
\midrule\noalign{}
\endhead
\bottomrule\noalign{}
\endlastfoot
Broadcast & Some endpoint accepted the transaction for propagation &
Inclusion, execution, or finality \\
Included & A block contains the transaction & Successful execution or
durable settlement \\
Receipt success & Top-level execution did not revert & Economic safety
or user intent \\
Event emitted & Code emitted a log & State changed as the event name
suggests \\
State updated & Canonical state reflects the change & The change was
safe or final enough for downstream use \\
Finalized/settled & Reversion is outside the accepted risk model &
Off-chain interpretation, indexer correctness, or user intent \\
\end{longtable}
}

For example:

\begin{lstlisting}
The transaction was broadcast.
The transaction was included.
The transaction succeeded with status 1.
The transaction emitted event X.
The transaction updated storage slot Y.
The transaction is finalized.
The transaction was observed by this indexer.
\end{lstlisting}

Precision matters because attackers exploit ambiguity.

\subsection{Receipts Summarize Execution Outcomes}

Ethereum transaction receipts include information such as block hash,
block number, gas used, logs, contract address for contract creation,
transaction index, type, and either a status or a pre-Byzantium state
root. Ethereum JSON-RPC documentation specifies that
\passthrough{\lstinline!status!} is \passthrough{\lstinline!1!} for
success or \passthrough{\lstinline!0!} for failure, and receipts include
logs and gas used fields.\footnote{ethereum.org, ``JSON-RPC API'',
  https://ethereum.org/developers/docs/apis/json-rpc/.}

Receipts are useful because they summarize execution without requiring
every observer to replay the transaction immediately.

But receipts are not a substitute for understanding state:

\begin{itemize}
\tightlist
\item
  A successful status means execution did not revert at the top level.
\item
  It does not mean the business operation was economically safe.
\item
  Logs may be emitted by contracts, but they do not prove that an
  off-chain interpretation is correct.
\item
  A receipt can disappear from canonical history if its block is reorged
  out before finality.
\item
  Indexers can misinterpret receipts or miss reorgs.
\end{itemize}

Receipt analysis asks:

\begin{lstlisting}
Was the transaction included?
Did it succeed or fail?
How much gas was used?
Which logs were emitted?
Which contract emitted each log?
Which state change do the logs claim to represent?
Is the block final enough for this application?
\end{lstlisting}

\subsection{Logs and Events Are Communication, Not Authority}

Solidity documentation describes events as an abstraction over the EVM
logging functionality. Event arguments are stored in transaction logs
associated with the emitting contract address; applications can
subscribe through client RPC interfaces.\footnote{Solidity
  Documentation, ``Events'',
  https://docs.soliditylang.org/en/latest/contracts.html\#events.}

Events are extremely useful for:

\begin{itemize}
\tightlist
\item
  Indexers.
\item
  UIs.
\item
  Analytics.
\item
  Off-chain automation.
\item
  Bridges and message watchers.
\item
  Monitoring.
\item
  Human debugging.
\end{itemize}

But events are not state accessible to contracts that emitted them, and
contracts cannot generally use their own past logs as authoritative
state. Solidity documentation also notes that event data must be
interpreted with knowledge of the event type, and anonymous events can
complicate interpretation.\footnote{Solidity Documentation, ``Events'',
  https://docs.soliditylang.org/en/latest/contracts.html\#events.}

Security consequences:

\begin{itemize}
\tightlist
\item
  An event can be emitted without the expected state change if the code
  is wrong.
\item
  A state change can occur without an expected event if the code omits
  it.
\item
  An indexer can treat an event as truth and become inconsistent with
  state.
\item
  A bridge that watches events must account for finality and reorgs.
\item
  Anonymous or misleading events can confuse tooling.
\item
  Off-chain systems must verify the emitting contract and event
  semantics.
\end{itemize}

Rule:

\begin{lstlisting}
Events are evidence.
State is authority.
\end{lstlisting}

Sometimes a protocol intentionally uses logs as part of a cross-domain
commitment path. Then the security model must say exactly which log,
from which contract, in which finalized block, under which proof or
watcher assumption, is accepted.

{\def\LTcaptype{none} % do not increment counter
\begin{longtable}[]{@{}
  >{\raggedright\arraybackslash}p{(\linewidth - 6\tabcolsep) * \real{0.2500}}
  >{\raggedright\arraybackslash}p{(\linewidth - 6\tabcolsep) * \real{0.2500}}
  >{\raggedright\arraybackslash}p{(\linewidth - 6\tabcolsep) * \real{0.2500}}
  >{\raggedright\arraybackslash}p{(\linewidth - 6\tabcolsep) * \real{0.2500}}@{}}
\toprule\noalign{}
\begin{minipage}[b]{\linewidth}\raggedright
Evidence type
\end{minipage} & \begin{minipage}[b]{\linewidth}\raggedright
Produced by
\end{minipage} & \begin{minipage}[b]{\linewidth}\raggedright
Useful for
\end{minipage} & \begin{minipage}[b]{\linewidth}\raggedright
Security warning
\end{minipage} \\
\midrule\noalign{}
\endhead
\bottomrule\noalign{}
\endlastfoot
State & Protocol/client execution & Authority for future validity & Must
be read from the correct canonical state \\
Receipt & Client execution result & Execution status, gas, logs, block
context & Can be reorged out or misread \\
Event/log & Contract or program execution & Off-chain discovery and
monitoring & Communicates; does not define authority by itself \\
Indexer view & Off-chain derived database & Search, UI, analytics,
automation & Can lag, omit reorgs, or interpret incorrectly \\
Simulation & Local execution against assumed state & Preflight and UX
warning & Not a promise about eventual inclusion or ordering \\
\end{longtable}
}

\subsection{Return Data and Revert Data Matter}

Execution can produce return data or revert data.

In EVM systems:

\begin{itemize}
\tightlist
\item
  A call may return bytes.
\item
  A call may revert with error data.
\item
  A low-level call may report success/failure without automatically
  bubbling the revert.
\item
  A token may return \passthrough{\lstinline!false!}, return no value,
  revert, or behave nonstandardly.
\item
  A wrapper library may interpret return behavior differently.
\end{itemize}

Solidity documentation states that \passthrough{\lstinline!revert!} and
failed \passthrough{\lstinline!require!} calls revert changes in the
current call and provide error data.\footnote{Solidity Documentation,
  ``Error handling: Assert, Require, Revert and Exceptions'',
  https://docs.soliditylang.org/en/latest/control-structures.html\#error-handling-assert-require-revert-and-exceptions.}

Security consequences:

\begin{itemize}
\tightlist
\item
  Ignoring return values can make a failed operation appear successful.
\item
  Assuming all tokens return booleans can break integrations.
\item
  Catching reverts can create partial-success semantics.
\item
  Low-level calls can hide failure if the result is not checked.
\item
  A transaction can succeed while an internal operation failed and was
  ignored.
\end{itemize}

Ask:

\begin{lstlisting}
Which call's success matters?
Is the return value checked?
Are errors bubbled, swallowed, or transformed?
Can partial progress remain after a failure?
Does the chain's atomicity rule cover this whole operation or only part of it?
\end{lstlisting}

\subsection{Atomicity Is Model-Specific}

People often say blockchain transactions are atomic. That is sometimes
true, but only at the correct boundary.

Examples:

\begin{itemize}
\tightlist
\item
  In Ethereum, if a top-level transaction reverts, state changes from
  that execution are reverted, but gas is still consumed. Internal calls
  can fail and be handled without reverting the top-level transaction.
\item
  In Solana, if any instruction in a transaction fails, the entire
  transaction fails and state changes are reverted, while fees are still
  charged.\footnote{Solana Foundation, ``Transactions'',
    https://solana.com/docs/core/transactions.}
\item
  In Cosmos SDK, failed state changes during delivery are reverted, but
  fee and ante-handler behavior must be understood in
  context.\footnote{Cosmos SDK v0.47, ``Transaction Lifecycle'',
    https://docs.cosmos.network/sdk/v0.47/learn/beginner/tx-lifecycle.}
\item
  In Bitcoin, a transaction either becomes part of canonical history or
  does not, but protocols built from multiple transactions do not get
  automatic atomicity across those transactions.
\item
  Cross-chain operations are almost never atomic in the same sense as a
  single-chain transaction.
\end{itemize}

Atomicity analysis asks:

\begin{lstlisting}
What exactly reverts?
What does not revert?
Are fees charged?
Are logs preserved?
Are nonce or sequence effects preserved?
Can internal failures be caught?
Does a multi-step protocol assume atomicity across multiple transactions or chains?
\end{lstlisting}

\subsection{Simulation Is Not Execution}

Wallets, bots, keepers, searchers, and protocols often simulate
transactions before submitting them.

Simulation is useful, but it is not the same as execution because:

\begin{itemize}
\tightlist
\item
  The state may change before inclusion.
\item
  The transaction may be ordered differently.
\item
  The block context may differ.
\item
  The account nonce may change.
\item
  Gas prices may change.
\item
  Oracle prices may update.
\item
  Another transaction may consume the same opportunity.
\item
  A different node or RPC may simulate against a different state view.
\item
  The transaction may land on another fork or not land at all.
\end{itemize}

Simulation answers:

\begin{lstlisting}
What would happen if this transaction ran now under this assumed state and context?
\end{lstlisting}

It does not answer:

\begin{lstlisting}
What will happen when this transaction is actually included?
\end{lstlisting}

This gap is central to MEV, liquidations, wallet UX, and transaction
risk.

\section{When a Valid Transaction Still Produces an Unsafe State}

\subsection{Validity Is Layered}

A transaction can pass many validity layers:

\begin{lstlisting}
well-encoded
properly signed
nonce or sequence correct
fees sufficient
accepted by a node
included in a block
executed without revert
receipt status success
finalized
\end{lstlisting}

And still be unsafe.

Why? Because these layers prove that the transaction satisfied explicit
machine rules. They do not prove that the protocol encoded the right
economic, authorization, or liveness properties.

Examples:

\begin{itemize}
\tightlist
\item
  A lending market accepts a borrow against a manipulated oracle price.
\item
  A vault accepts a deposit that mints zero shares because of rounding.
\item
  A bridge accepts a message before the source event is final.
\item
  A governance proposal executes malicious calls after satisfying
  quorum.
\item
  A user signs an unlimited approval through a malicious frontend.
\item
  A liquidation succeeds because the transaction was ordered after a
  temporary price update.
\item
  A keeper transaction succeeds but pays an attacker-controlled address
  because configuration was stale.
\end{itemize}

The transaction is not invalid. The protocol design is insufficient.

\subsection{Successful Execution Can Encode a Failed Assumption}

A successful transaction may mean:

\begin{lstlisting}
The code path completed.
\end{lstlisting}

It does not necessarily mean:

\begin{lstlisting}
The oracle price was fair.
The user understood the signature.
The admin action was legitimate.
The bridge source state was final.
The governance outcome was socially acceptable.
The vault shares were fairly priced.
The recipient can use the received asset.
The protocol remains solvent.
\end{lstlisting}

A receipt status can tell you whether the execution reverted. It cannot
tell you whether the transition preserved intent.

Security review must therefore recover the implied assumptions behind a
transaction:

\begin{lstlisting}
This transition is safe if:
  price is fresh and manipulation-resistant
  signer intended this operation
  nonce is unique and domain-scoped
  output is final enough
  role holder is trusted or constrained
  resource limits leave the exit path live
\end{lstlisting}

If the ``if'' clause is false, successful execution can be the exploit
path.

\subsection{Unsafe State From Correct Ordering Rules}

Ordering is part of transaction semantics.

Two valid transactions can produce different outcomes depending on
order:

\begin{lstlisting}
tx_A: update oracle price
tx_B: borrow against collateral
\end{lstlisting}

If \passthrough{\lstinline!tx\_B!} executes before
\passthrough{\lstinline!tx\_A!}, the borrower sees the old price. If
\passthrough{\lstinline!tx\_B!} executes after
\passthrough{\lstinline!tx\_A!}, the borrower sees the new price. Both
orderings may be valid under chain rules. Only one may be safe under the
application's assumptions.

Ordering risk appears in:

\begin{itemize}
\tightlist
\item
  Front-running.
\item
  Back-running.
\item
  Sandwiching.
\item
  Liquidations.
\item
  Oracle updates.
\item
  NFT mints.
\item
  Governance proposal execution.
\item
  Bridge message delivery.
\item
  Withdrawals from queues.
\item
  Auctions.
\end{itemize}

Do not ask only:

\begin{lstlisting}
Can this transaction execute?
\end{lstlisting}

Ask:

\begin{lstlisting}
What can execute before it?
What can execute after it?
Who can observe it while pending?
Who can choose or influence ordering?
What state assumptions become false if ordering changes?
\end{lstlisting}

\subsection{Unsafe State From Partial Mental Models}

Many transaction bugs come from developers assuming the transaction has
only one effect.

Example:

\begin{lstlisting}
transfer token from user to vault
mint shares to user
\end{lstlisting}

The developer thinks there are two effects. A real transaction might
include:

\begin{itemize}
\tightlist
\item
  Token hook calls.
\item
  Fee-on-transfer behavior.
\item
  ERC-777-style callbacks.
\item
  Reentrant calls.
\item
  Oracle updates.
\item
  Share price rounding.
\item
  Event emission.
\item
  Internal accounting.
\item
  External strategy deposit.
\item
  Gas griefing.
\item
  Failure behavior from a nonstandard token.
\end{itemize}

Security review expands the mental model until every relevant effect is
visible.

\subsection{Unsafe State From Stale Context}

A transaction is constructed at one time and executed later.

Between construction and execution:

\begin{itemize}
\tightlist
\item
  Prices can move.
\item
  Nonces can change.
\item
  Allowances can change.
\item
  Blocks can advance.
\item
  Governance proposals can queue.
\item
  Liquidation status can change.
\item
  A pool can be manipulated.
\item
  A target contract can be upgraded.
\item
  A bridge source event can be reorged.
\item
  A Solana recent blockhash can expire.
\item
  A Sui epoch can advance beyond expiration.
\end{itemize}

A transaction that was safe when signed may be unsafe when executed.

Defenses include:

\begin{itemize}
\tightlist
\item
  Slippage bounds.
\item
  Deadlines.
\item
  Nonces.
\item
  Domain separation.
\item
  Expected state checks.
\item
  Minimum output amounts.
\item
  Freshness checks.
\item
  Finality delays.
\item
  Commit-reveal mechanisms.
\item
  Explicit target and version binding.
\end{itemize}

Every transaction that depends on mutable context should say:

\begin{lstlisting}
Execute only if these assumptions are still true.
\end{lstlisting}

\subsection{Unsafe State From Successful Failure}

Sometimes attackers want a transaction to fail.

Examples:

\begin{itemize}
\tightlist
\item
  Griefing by causing another user's transaction to revert.
\item
  Filling blocks or queues with failing transactions.
\item
  Forcing a keeper to pay fees repeatedly.
\item
  Causing a liquidation to fail so bad debt grows.
\item
  Making a withdrawal batch fail because one recipient reverts.
\item
  Locking accounts or objects so competitors cannot execute.
\item
  Triggering expensive failure paths.
\end{itemize}

Failure is an outcome. It can be profitable.

Security review should ask:

\begin{lstlisting}
Who pays when this fails?
Who benefits when it fails?
Can one failing operation block a batch?
Can failure leave state partially updated?
Can repeated failure exhaust resources?
\end{lstlisting}

\subsection{Transaction Safety Checklist}

For any important transaction path, review:

\begin{lstlisting}
Authorization:
  Who signed or authorized?
  Is the signature scoped to the correct chain, contract, function, account, nonce, and deadline?

State:
  Which state is read?
  Which state is written?
  Which assumed state can change before inclusion?

Ordering:
  Can the transaction be observed, copied, frontrun, backrun, delayed, or censored?

Resources:
  What gas, fees, compute units, size limits, or account limits apply?
  Can the transition remain executable under worst-case state?

Failure:
  What happens if execution reverts or partially fails?
  Who pays?
  Are nonce, fee, log, or queue effects preserved?

Results:
  Which receipt, event, log, or state update proves success?
  Is the application relying on logs instead of state?

Finality:
  When is the result safe to depend on?
  What happens if the block or source event is reorged?

Intent:
  What did the user or protocol intend?
  Which invariant should the transaction preserve?
\end{lstlisting}

If this checklist feels long, that is because transactions are where
many separate assumptions collide.

\section{Worked Example: A Swap Transaction}

Consider a user signing a transaction through a decentralized exchange
frontend:

\begin{lstlisting}
swapExactTokensForTokens(
  tokenIn = USDC,
  tokenOut = WETH,
  amountIn = 10,000 USDC,
  minAmountOut = 3.8 WETH,
  path = [USDC, WETH],
  recipient = user,
  deadline = now + 10 minutes
)
\end{lstlisting}

The frontend says:

\begin{lstlisting}
Swap 10,000 USDC for at least 3.8 WETH.
\end{lstlisting}

A transaction review asks:

\begin{lstlisting}
Authorization:
  Did the user sign the swap transaction?
  Did the router already have allowance?
  Was there a separate approval transaction?

State:
  Which pool reserves determine the price?
  Is the token standard or nonstandard?
  Does either token take fees on transfer?

Ordering:
  Can the transaction be sandwiched?
  Can a searcher move the pool price before and after the swap?

Resources:
  Is gas limit sufficient?
  Can the transaction be stuck because fee settings are too low?

Failure:
  If output is below minAmountOut, does it revert?
  If it reverts, what fees are consumed?
  Is the approval still outstanding?

Results:
  Did the transaction status succeed?
  Did balances change as expected?
  Are events consistent with state?

Finality:
  Is the block final enough for the UI or downstream protocol to rely on?

Intent:
  Did the user intend this exact router, path, recipient, amount, and deadline?
\end{lstlisting}

The most important security control in this transaction is not that the
user signed it. The signature only proves authorization. The safety
controls are the bound on \passthrough{\lstinline!minAmountOut!}, the
deadline, the target contract, the token assumptions, and the user's
ability to understand what they are signing.

The same transaction as a trace:

{\def\LTcaptype{none} % do not increment counter
\begin{longtable}[]{@{}
  >{\raggedright\arraybackslash}p{(\linewidth - 4\tabcolsep) * \real{0.3333}}
  >{\raggedright\arraybackslash}p{(\linewidth - 4\tabcolsep) * \real{0.3333}}
  >{\raggedright\arraybackslash}p{(\linewidth - 4\tabcolsep) * \real{0.3333}}@{}}
\toprule\noalign{}
\begin{minipage}[b]{\linewidth}\raggedright
Step
\end{minipage} & \begin{minipage}[b]{\linewidth}\raggedright
Evidence
\end{minipage} & \begin{minipage}[b]{\linewidth}\raggedright
Security question
\end{minipage} \\
\midrule\noalign{}
\endhead
\bottomrule\noalign{}
\endlastfoot
User signs & Signature over router call & Did the signed payload bind
the intended router, path, recipient, amount, deadline, and chain? \\
Transaction waits & Mempool/private route observation & Who can see it,
copy it, delay it, or order around it? \\
Router executes & Pool reserves and token behavior & Are price,
slippage, transfer fees, and callbacks within assumptions? \\
Receipt appears & Status, gas used, emitted logs & Did execution
succeed, and do logs match actual state changes? \\
Balances change & Token balances and pool reserves & Was the user paid
at least \passthrough{\lstinline!minAmountOut!}, and did allowances
remain risky? \\
Block settles & Confirmation/finality signal & Is the result durable
enough for the UI or downstream protocol? \\
\end{longtable}
}

Now consider a malicious variant:

\begin{lstlisting}
approve(attacker, type(uint256).max)
\end{lstlisting}

This transaction may:

\begin{itemize}
\tightlist
\item
  Be well-encoded.
\item
  Be signed by the user.
\item
  Pay enough gas.
\item
  Be included in a valid block.
\item
  Execute successfully.
\item
  Emit an \passthrough{\lstinline!Approval!} event.
\end{itemize}

It is still unsafe for the user. The failure is not signature
validation. The failure is that the user-side authority was tricked into
authorizing a harmful state transition.

\section{Practical Discipline: Transaction Review Notes}

When reviewing a transaction path, write notes in this format:

\begin{lstlisting}
Transaction:
  Name or entry point.

Intent:
  What the user/protocol thinks this does.

Signed authority:
  Who signs, or which account/role authorizes.

Replay scope:
  Nonce, sequence, object version, blockhash, chain ID, domain, deadline.

State touched:
  Reads and writes.

External assumptions:
  Oracles, token behavior, bridge finality, keeper behavior, frontend correctness.

Resource assumptions:
  Gas, compute, fees, size, account count, block limits.

Ordering assumptions:
  What must not happen before/after.

Failure behavior:
  What reverts, what persists, who pays.

Observable result:
  Receipt status, logs, state changes, confirmations/finality.

Invariant:
  What must remain true afterward.
\end{lstlisting}

This format forces transaction analysis to connect mechanics to security
properties.

\section{Review Questions}

\begin{enumerate}
\def\labelenumi{\arabic{enumi}.}
\tightlist
\item
  What is the difference between a transaction, an instruction, a
  message, and a state transition?
\item
  In a UTXO model, what state is consumed and what state is created?
\item
  Why can a zero-ETH Ethereum transaction still move valuable assets?
\item
  Why do Solana account lists and signer/writable flags matter for
  security?
\item
  What does a transaction signature prove, and what does it not prove?
\item
  How do nonces, sequence numbers, blockhashes, chain IDs, and deadlines
  constrain replay?
\item
  What is the difference between inclusion, execution success, and
  settlement?
\item
  Why can events, receipts, and simulations mislead off-chain systems if
  treated as authority?
\end{enumerate}

\section{Applied Exercises}

\subsection{Exercise 1: Transaction Anatomy and Evidence}

Compare three transaction models:

\begin{itemize}
\tightlist
\item
  Bitcoin transaction spending a UTXO.
\item
  Ethereum transaction calling an ERC-20
  \passthrough{\lstinline!transfer!}.
\item
  Solana transaction with two instructions.
\end{itemize}

For each, identify:

\begin{itemize}
\tightlist
\item
  Authorization.
\item
  Replay-prevention mechanism.
\item
  State consumed or written.
\item
  Resource limits or fees.
\item
  Failure behavior.
\item
  Evidence an observer should check to confirm success.
\end{itemize}

Then choose one real or hypothetical transaction receipt/result and
identify:

\begin{itemize}
\tightlist
\item
  Signer or authority.
\item
  Replay-prevention mechanism.
\item
  State touched.
\item
  Logs or events.
\item
  Revert or failure behavior.
\item
  Settlement assumption.
\item
  One way the transaction could be valid but unsafe.
\end{itemize}

\subsection{Exercise 2: Replay Scope}

A protocol accepts off-chain signatures authorizing withdrawals:

\begin{lstlisting}
withdraw(user, amount, signature)
\end{lstlisting}

The signed message includes only:

\begin{lstlisting}
user
amount
\end{lstlisting}

Answer:

\begin{itemize}
\tightlist
\item
  Which replay attacks are possible?
\item
  Which fields should be added to the signed message?
\item
  Which state should be updated when the signature is consumed?
\item
  Should the signature expire?
\item
  Should the signature be valid across contract upgrades?
\end{itemize}

\subsection{Exercise 3: Event vs State}

A bridge watcher sees this event:

\begin{lstlisting}
Deposit(address indexed user, uint256 amount, uint256 nonce)
\end{lstlisting}

The watcher mints wrapped tokens on another chain.

Answer:

\begin{itemize}
\tightlist
\item
  What must the watcher verify besides the event data?
\item
  Why is the emitting contract address important?
\item
  What finality assumption is required?
\item
  What replay-protection state should exist on the destination chain?
\item
  What happens if the source block is reorged?
\end{itemize}

\subsection{Exercise 4: Resource Limits and Valid-But-Unsafe Outcomes}

A protocol has a function:

\begin{lstlisting}
processAllWithdrawals()
\end{lstlisting}

It loops over every pending withdrawal and pays each recipient.

Answer:

\begin{itemize}
\tightlist
\item
  What resource limit can break this function?
\item
  How can an attacker make the problem worse?
\item
  Is this a safety failure, liveness failure, or both?
\item
  What design patterns would reduce the risk?
\end{itemize}

Then explain why each transaction may be valid and still unsafe:

\begin{enumerate}
\def\labelenumi{\arabic{enumi}.}
\tightlist
\item
  A user signs an unlimited token approval to a malicious spender.
\item
  A lending market accepts a borrow after a one-block oracle
  manipulation.
\item
  A governance proposal executes after satisfying quorum, but transfers
  the treasury to the proposer.
\item
  A bridge releases funds after seeing a source-chain event with too few
  confirmations.
\item
  A liquidation transaction succeeds because it was ordered immediately
  after a temporary price update.
\end{enumerate}

\section{Key Takeaways}

\begin{itemize}
\tightlist
\item
  Transactions are signed, metered, ordered requests to execute state
  transitions.
\item
  Transaction anatomy is a security interface: every field determines
  authority, replay scope, resource use, or execution context.
\item
  UTXO, account, instruction-based, object-based, and Cosmos-style
  systems package transactions differently.
\item
  Signatures prove authorization over specific bytes and domains. They
  do not prove user understanding or protocol safety.
\item
  Replay prevention must be scoped to the authority and asset being
  consumed.
\item
  Fees, gas, compute budgets, size limits, and block limits are security
  mechanisms.
\item
  Inclusion, execution success, receipt status, event emission, and
  finality are different facts.
\item
  Logs and events are evidence for observers, not a substitute for
  state.
\item
  Simulation is useful but not execution.
\item
  A transaction can be valid, successful, finalized, and still unsafe if
  the protocol encoded the wrong assumptions.
\end{itemize}

\section{Further Reading}

\begin{itemize}
\tightlist
\item
  ethereum.org,
  \href{https://ethereum.org/developers/docs/transactions/}{``Transactions''}.
  Practical overview of Ethereum transaction fields, signing, calldata,
  gas fields, and lifecycle.
\item
  ethereum.org, \href{https://ethereum.org/developers/docs/gas/}{``Gas
  and Fees''}. Technical overview of Ethereum gas, gas limits, base fee,
  priority fee, and transaction costs.
\item
  Vitalik Buterin,
  \href{https://eips.ethereum.org/EIPS/eip-155}{``EIP-155: Simple Replay
  Attack Protection''}. The Ethereum replay-protection change that binds
  transaction signing to chain ID.
\item
  Vitalik Buterin et al.,
  \href{https://eips.ethereum.org/EIPS/eip-1559}{``EIP-1559: Fee Market
  Change for ETH 1.0 Chain''}. The fee-market change introducing base
  fee mechanics and dynamic fee transactions.
\item
  Bitcoin Developer Guide,
  \href{https://developer.bitcoin.org/devguide/transactions.html}{``Transactions''}.
  Detailed explanation of Bitcoin transaction inputs, outputs, scripts,
  and validation.
\item
  Solana Foundation,
  \href{https://solana.com/docs/core/transactions}{``Transactions''}.
  Reference for Solana transaction structure, signatures, atomicity,
  blockhash expiry, and limits.
\item
  Solana Foundation, \href{https://solana.com/docs/core/fees}{``Fees''}.
  Reference for base fees, prioritization fees, and
  compute-budget-related limits.
\item
  Cosmos SDK v0.47,
  \href{https://docs.cosmos.network/sdk/v0.47/learn/beginner/tx-lifecycle}{``Transaction
  Lifecycle''}. Versioned reference for transaction creation, mempool
  checks, ante handling, delivery, and commit.
\item
  Sui Foundation,
  \href{https://docs.sui.io/develop/transactions/txn-overview}{``Transaction
  Overview''}. Overview of transaction metadata, object inputs and
  outputs, gas input, expiration, and effects.
\item
  ethereum.org,
  \href{https://ethereum.org/developers/docs/apis/json-rpc/}{``JSON-RPC
  API''}. Useful for transaction receipts, logs, status, and execution
  client interaction.
\item
  Solidity Documentation,
  \href{https://docs.soliditylang.org/en/latest/contracts.html\#events}{``Events''}.
  Reference for events, logs, topics, and event interpretation.
\item
  Solidity Documentation,
  \href{https://docs.soliditylang.org/en/latest/control-structures.html\#error-handling-assert-require-revert-and-exceptions}{``Error
  handling: Assert, Require, Revert and Exceptions''}. Reference for
  \passthrough{\lstinline!require!}, \passthrough{\lstinline!revert!},
  exceptions, and state rollback.
\end{itemize}
